\chapter{المستقيمات والدوائر والزوايا}\label{ch:g6:lines}

تبدأ الهندسة بمسطرة وكوس وفرجار. ويثبّت هذا الفصل
المفردات --- النقط \omterm{def:g6:lines:objects}{والقطع} وأنصاف \omterm{def:g6:lines:objects}{المستقيمات} \omterm{def:g6:lines:objects}{والمستقيمات} والدوائر ---
ويقدّم الوضعين الخاصّين \omterm{def:g6:lines:objects}{للمستقيمين} (التوازي
والتعامد)، ويعلّم كيف تُقاس الزوايا وتُرسم.

\section{النقط والقطع والمستقيمات}

\begin{definition}[الأشياء الأساسية]\label{def:g6:lines:objects}
يمرّ بنقطتين متمايزتين $A$ و $B$:
\begin{itemize}
  \item \emph{القطعة}\index{قطعة} $[AB]$: وهي جزء المستقيم
        المحصور بين $A$ و $B$ (ولها طول يُكتب $AB$)؛
  \item و\emph{نصف المستقيم} $[AB)$: يبدأ من $A$، ويمرّ بالنقطة $B$، ثم
        يمتدّ بلا نهاية؛
  \item و\emph{المستقيم}\index{مستقيم} $(AB)$: يمتدّ بلا نهاية في
        الجهتين. والنقطتان تعيّنان مستقيمًا واحدًا لا غير.
\end{itemize}
والنقط الواقعة على مستقيم واحد تُسمّى \emph{مستقيمية}.
\end{definition}

\begin{omfigure}
\begin{tikzpicture}[scale=1.0]
  \fill (0,1.8) circle (1.5pt) node[above, font=\small] {$A$};
  \fill (2.5,1.8) circle (1.5pt) node[above, font=\small] {$B$};
  \draw[thick, omDef] (0,1.8) -- (2.5,1.8);
  \node[right, font=\small] at (2.9,1.8) {القطعة $[AB]$};
  \fill (0,0.9) circle (1.5pt) node[above, font=\small] {$A$};
  \fill (2.5,0.9) circle (1.5pt) node[above, font=\small] {$B$};
  \draw[thick, omThm] (0,0.9) -- (4.1,0.9);
  \node[right, font=\small] at (4.3,0.9) {نصف المستقيم $[AB)$};
  \fill (0,0) circle (1.5pt) node[above, font=\small] {$A$};
  \fill (2.5,0) circle (1.5pt) node[above, font=\small] {$B$};
  \draw[thick, omMeth] (-1.4,0) -- (4.1,0);
  \node[right, font=\small] at (4.3,0) {المستقيم $(AB)$};
\end{tikzpicture}

{\small النقطتان نفسهما، وثلاثة أشياء مختلفة. والقوسان يقولان ما
يتوقّف وما يستمرّ: $[$ يتوقّف، و $($ يستمرّ.}
\end{omfigure}

\section{المستقيمات المتوازية والمتعامدة}

\begin{definition}[المتعامدان والمتوازيان]\label{def:g6:lines:perp}
يكون \omterm{def:g6:lines:objects}{المستقيمان} \emph{متعامدين}\index{مستقيمان متعامدان} حين
يتقاطعان في \omterm{def:g3:shapes:rightangle}{زاوية قائمة}؛ ويُكتب $d_1 \perp d_2$، وتُعلَّم \omterm{def:g3:shapes:rightangle}{الزاوية
القائمة} بمربع صغير. ويكونان
\emph{متوازيين}\index{مستقيمان متوازيان} حين لا يتقاطعان أبدًا، مهما
اِمتدّا؛ ويُكتب $d_1 \parallel d_2$.
\end{definition}

\begin{omfigure}
\begin{tikzpicture}[scale=1.0]
  \draw[thick] (-1.8,0) -- (1.8,0) node[right, font=\small] {$d_1$};
  \draw[thick, omDef] (0.6,-1.2) -- (0.6,1.4)
    node[above, font=\small] {$d_2$};
  \draw[thin] (0.6,0) ++(-0.25,0) -- ++(0,0.25) -- ++(0.25,0);
  \node[below, font=\small] at (0,-1.5) {$d_1 \perp d_2$};
  \begin{scope}[xshift=6.5cm]
  \draw[thick] (-1.8,-0.5) -- (1.8,0.3) node[right, font=\small] {$d_3$};
  \draw[thick, omThm] (-1.8,-1.3) -- (1.8,-0.5)
    node[right, font=\small] {$d_4$};
  \node[below, font=\small] at (0,-1.5) {$d_3 \parallel d_4$};
  \end{scope}
\end{tikzpicture}

{\small تقاطع في \omterm{def:g3:shapes:rightangle}{زاوية قائمة} (مُعلَّمة بالمربع الصغير)، \omterm{def:g6:lines:perp}{ومستقيمان
متوازيان}: الاتجاه نفسه، ولا نقطة تقاطع.}
\end{omfigure}

\begin{proposition}[واقعتان مفيدتان]\label{prop:g6:lines:facts}
\begin{enumerate}
  \item إذا كان \omterm{def:g6:lines:objects}{مستقيمان} عموديين على \omterm{def:g6:lines:objects}{مستقيم} ثالث، فهما
        متوازيان.
  \item وإذا كان \omterm{def:g6:lines:objects}{مستقيمان} \omterm{def:g6:lines:perp}{متوازيين}، فكل \omterm{def:g6:lines:objects}{مستقيم} عمودي على أحدهما
        عمودي على الآخر.
\end{enumerate}
\end{proposition}
\admitted

\begin{method}[الرسم بالكوس]\label{met:g6:lines:setsquare}
لرسم العمودي على \omterm{def:g6:lines:objects}{مستقيم} $d$ المارّ بنقطة $P$:
\begin{enumerate}
  \item ضع أحد حرفَي \omterm{def:g3:shapes:rightangle}{الزاوية القائمة} من الكوس على $d$؛
  \item أزلق الكوس على طول $d$ حتى يبلغ حرفه الآخر النقطة
        $P$؛
  \item اُرسم \omterm{def:g6:lines:objects}{المستقيم} على طول ذلك الحرف، وعلّم \omterm{def:g3:shapes:rightangle}{الزاوية القائمة}.
\end{enumerate}
وللحصول على \omterm{def:g6:lines:perp}{الموازي} المارّ بالنقطة $P$: اُرسم عموديًّا على $d$، ثم
العمودي على \emph{ذلك} \omterm{def:g6:lines:objects}{المستقيم} المارّ بالنقطة $P$
(\cref{prop:g6:lines:facts}).
\end{method}

\section{الدوائر}

\begin{definition}[الدائرة]\label{def:g6:lines:circle}
\emph{الدائرة}\index{دائرة} التي مركزها $O$ ونصف قطرها $r$ هي مجموعة
كل النقط التي تبعد المسافة $r$ بالضبط عن $O$. \omterm{def:g6:lines:objects}{والقطعة} الواصلة من
المركز إلى الدائرة \emph{\omterm{def:g3:shapes:shapes}{نصف قطر}}؛ \omterm{def:g6:lines:objects}{والقطعة} الواصلة بين نقطتين من
الدائرة مارّةً بالمركز \emph{\omterm{def:g4:geometry:circle}{قطر}} --- وطوله
$2r$؛ \omterm{def:g6:lines:objects}{والقطعة} الواصلة بين نقطتين من الدائرة \emph{وتر}.
\end{definition}

\begin{omfigure}
\begin{tikzpicture}[scale=1.15]
  \draw[thick] (0,0) circle (1.5);
  \fill (0,0) circle (1.5pt) node[below left, font=\small] {$O$};
  \fill (1.5,0) circle (1.5pt) node[right, font=\small] {$M$};
  \draw[omDef, thick] (0,0) -- (1.5,0)
    node[midway, above, font=\small] {نصف القطر};
  \fill (-1.06,1.06) circle (1.5pt) node[above left, font=\small] {$A$};
  \fill (1.06,-1.06) circle (1.5pt) node[below right, font=\small] {$B$};
  \draw[omThm, thick] (-1.06,1.06) -- (1.06,-1.06)
    node[pos=0.25, right=2pt, font=\small] {القطر};
  \fill (0.52,1.41) circle (1.5pt) node[above, font=\small] {$C$};
  \fill (1.41,0.52) circle (1.5pt) node[right, font=\small] {$D$};
  \draw[omMeth, thick] (0.52,1.41) -- (1.41,0.52)
    node[midway, above right=-2pt, font=\small] {وتر};
\end{tikzpicture}

{\small دائرة فيها \omterm{def:g3:shapes:shapes}{نصف قطر} $[OM]$، \omterm{def:g4:geometry:circle}{وقطر} $[AB]$ (وهو وتر
يمرّ بالمركز، وطوله ضعف \omterm{def:g3:shapes:shapes}{نصف القطر})، ووتر $[CD]$.}
\end{omfigure}

\begin{example}\label{ex:g6:lines:circle}
«اُرسم الدائرة التي مركزها $O$ وتمرّ بالنقطة $A$»: اِفتح الفرجار
من $O$ إلى $A$ --- \omterm{def:g3:shapes:shapes}{فنصف القطر} هو المسافة $OA$ --- ثم أدِر. وكل
نقطة من هذه الدائرة تبعد عن $O$ مثل $A$.
\end{example}

\section{الزوايا}

\begin{definition}[الزاوية]\label{def:g6:lines:angle}
نصفا \omterm{def:g6:lines:objects}{المستقيمين} $[AB)$ و $[AC)$ اللذان لهما نقطة البداية نفسها $A$ يكوّنان
\emph{الزاوية}\index{زاوية} $\widehat{BAC}$؛ والنقطة $A$ هي
\emph{رأسها} (وهي دائمًا الحرف الأوسط!). وتُقاس الزوايا
\emph{بالدرجات} ($^\circ$)، من $0^\circ$ إلى $360^\circ$ للدورة
الكاملة. والزاوية:
\begin{itemize}
  \item \emph{قائمة} إذا كان قياسها $90^\circ$؛
  \item و\emph{حادّة} إذا كان قياسها أصغر من $90^\circ$؛
  \item و\emph{منفرجة} إذا كان قياسها بين $90^\circ$ و
        $180^\circ$؛
  \item و\emph{مستقيمة} إذا كان قياسها $180^\circ$ (فيكوّن نصفا
        \omterm{def:g6:lines:objects}{المستقيمين} مستقيمًا).
\end{itemize}
\end{definition}

\begin{omfigure}
\begin{tikzpicture}[scale=0.9]
  \draw[thick] (0,0) -- (1.8,0);
  \draw[thick] (0,0) -- (40:1.8);
  \draw[omDef, ->] (0.7,0) arc (0:40:0.7);
  \node[omDef, font=\small] at (20:1.0) {$40^\circ$};
  \node[below, font=\small] at (0.9,-0.3) {حادّة};
  \begin{scope}[xshift=3.6cm]
  \draw[thick] (0,0) -- (1.8,0);
  \draw[thick] (0,0) -- (0,1.8);
  \draw[thin] (0.3,0) -- (0.3,0.3) -- (0,0.3);
  \node[below, font=\small] at (0.9,-0.3) {قائمة};
  \end{scope}
  \begin{scope}[xshift=7.2cm]
  \draw[thick] (0,0) -- (1.8,0);
  \draw[thick] (0,0) -- (135:1.8);
  \draw[omThm, ->] (0.6,0) arc (0:135:0.6);
  \node[omThm, font=\small] at (65:1.0) {$135^\circ$};
  \node[below, font=\small] at (0.9,-0.3) {منفرجة};
  \end{scope}
  \begin{scope}[xshift=11.6cm]
  \draw[thick] (-1.5,0) -- (1.5,0);
  \fill (0,0) circle (1.2pt);
  \draw[omMeth, ->] (0.5,0) arc (0:180:0.5);
  \node[below, font=\small] at (0,-0.3) {مستقيمة};
  \end{scope}
\end{tikzpicture}

{\small عائلات الزوايا الأربع. \omterm{def:g3:shapes:rightangle}{والزاوية القائمة} ($90^\circ$) هي
المرجع: فالحادّة أصغر منها، والمنفرجة أكبر.}
\end{omfigure}

\begin{method}[قياس زاوية بالمنقلة]\label{met:g6:lines:protractor}
\begin{enumerate}
  \item ضع مركز المنقلة على رأس الزاوية
        بالضبط؛
  \item حاذِ خطّها $0^\circ$ مع أحد ضلعَي الزاوية؛
  \item اِقرأ التدريجة التي يقطعها الضلع الآخر --- مستعملًا
        السلّم الذي يبدأ من $0$ على الضلع المحاذي؛
  \item وتحقّق بالنظر: فالزاوية الحادّة يجب أن تُقرأ أصغر من
        $90^\circ$، والمنفرجة أكبر.
\end{enumerate}
\end{method}

\begin{example}\label{ex:g6:lines:estimate}
قبل القياس، قدِّر! فالزاوية المفتوحة أكثر قليلًا من ركن
ورقة تزيد قليلًا على $90^\circ$؛ ونصف \omterm{def:g3:shapes:rightangle}{الزاوية القائمة}
$45^\circ$؛ وثلثها $30^\circ$. وإن قالت منقلتك
$150^\circ$ لزاوية تبدو حادّة، فقد قرأت السلّم الخطأ:
والقياس الصحيح هو $180^\circ - 150^\circ = 30^\circ$.
\end{example}

\section{تمارين}

\begin{exercise}[$\star$]\label{exo:g6:lines:1}
اُرسم ثلاث نقط $A$ و $B$ و $C$ غير مستقيمية. واُرسم بألوان مختلفة:
\omterm{def:g6:lines:objects}{القطعة} $[AB]$، ونصف \omterm{def:g6:lines:objects}{المستقيم} $[CA)$، \omterm{def:g6:lines:objects}{والمستقيم} $(BC)$.
\end{exercise}

\begin{exercise}[$\star$]\label{exo:g6:lines:2}
صواب أم خطأ؟ «$[AB]$ و $[BA]$ هما \omterm{def:g6:lines:objects}{القطعة} نفسها.»
«$[AB)$ و $[BA)$ هما نصف \omterm{def:g6:lines:objects}{المستقيم} نفسه.» «$(AB)$ و $(BA)$ هما
\omterm{def:g6:lines:objects}{المستقيم} نفسه.» اِشرح كل جواب.
\end{exercise}

\begin{exercise}[$\star$]\label{exo:g6:lines:3}
اُرسم مستقيمًا $d$ ونقطة $P$ لا تنتمي إلى $d$. واِبنِ بالكوس:
العمودي على $d$ المارّ بالنقطة $P$، ثم \omterm{def:g6:lines:perp}{الموازي} \omterm{def:g6:lines:objects}{للمستقيم} $d$
المارّ بالنقطة $P$. وصِف خطواتك.
\end{exercise}

\begin{exercise}[$\star$]\label{exo:g6:lines:4}
\omterm{def:g6:lines:objects}{المستقيمان} $a$ و $b$ عموديان كلاهما على \omterm{def:g6:lines:objects}{مستقيم} $c$، \omterm{def:g6:lines:objects}{ومستقيم} رابع
$e$ عمودي على $a$. فماذا تقول عن $a$ و $b$؟
وعن $e$ و $c$؟ بَرِّر حسب \cref{prop:g6:lines:facts}.
\end{exercise}

\begin{exercise}[$\star$]\label{exo:g6:lines:5}
اُرسم دائرة مركزها $O$ ونصف قطرها $3$ سم. وضع نقطة $M$ على
الدائرة، ونقطة $N$ داخلها، ونقطة $P$ خارجها. فماذا تقول
عن المسافات $OM$ و $ON$ و $OP$ بالمقارنة مع $3$ سم؟
\end{exercise}

\begin{exercise}[$\star$]\label{exo:g6:lines:6}
\omterm{def:g4:geometry:circle}{قطر} دائرة $9$ سم. فما نصف قطرها؟ ودائرة أخرى نصف قطرها
$2.6$ سم: فما قطرها؟
\end{exercise}

\begin{exercise}[$\star$]\label{exo:g6:lines:7}
سمِّ الزاوية المعلَّمة بثلاثة \omterm{def:g5:solids:def}{أحرف}، ثم صنّفها (حادّة، قائمة،
منفرجة، مستقيمة): زاوية قياسها $72^\circ$ رأسها $R$ بين نصفَي \omterm{def:g6:lines:objects}{مستقيمين}
نحو $S$ و $T$؛ وزاوية قياسها $148^\circ$ رأسها $B$ بين نصفَي \omterm{def:g6:lines:objects}{مستقيمين}
نحو $A$ و $C$؛ وزاوية قياسها $90^\circ$ رأسها $O$ بين نصفَي \omterm{def:g6:lines:objects}{مستقيمين}
نحو $M$ و $N$.
\end{exercise}

\begin{exercise}[$\star$]\label{exo:g6:lines:8}
قدِّر، ثم قِس بالمنقلة، زوايا مثلث ترسمه بنفسك الثلاث.
واجمع القياسات الثلاثة: فماذا تجد؟
(اِحتفظ بجوابك من أجل \cref{ch:g7:triangles}.)
\end{exercise}

\begin{exercise}[$\star$]\label{exo:g6:lines:9}
اُرسم بالمنقلة زاوية $\widehat{xOy}$ قياسها $65^\circ$، ثم
زاوية قياسها $130^\circ$. فكيف تحصل على الثانية من الأولى
دون منقلة؟
\end{exercise}

\begin{exercise}[$\star\star$]\label{exo:g6:lines:10}
قريتان $A$ و $B$ مرسومتان على خريطة. فأين تقع النقط التي
تبعد $3$ سم عن $A$ \emph{و} $2$ سم عن $B$؟ اُرسم صورة
تبيّن كم نقطة كهذه يمكن أن توجد ($0$ أو $1$ أو $2$ بحسب
المسافة $AB$).
\end{exercise}

\begin{exercise}[$\star\star$]\label{exo:g6:lines:11}
تشير ساعة إلى 3 تمامًا: فما الزاوية بين العقربين؟ والسؤال نفسه
عند 5 تمامًا، وعند 6 تمامًا. (الدورة الكاملة $360^\circ$
في $12$ ساعة.)
\end{exercise}

\begin{exercise}[$\star\star\star$]\label{exo:g6:lines:12}
اُرسم \omterm{def:g6:lines:objects}{قطعة} $[AB]$ طولها $6$ سم. واِبنِ النقطة $C$ التي تحقّق
$AC = BC = 6$ سم، بالفرجار وحده، وقِس الزاوية
$\widehat{CAB}$. فأيّ مثلث بنيت، وماذا
تخمّن عن زواياه؟
\end{exercise}

\section{مسألة: هندسة وجه الساعة}

\begin{problem}[{مسألة نهاية الأسبوع --- الزوايا كسورًا من
الدورة: قراءتها على ساعة، وتصيّد اللحظات التي يلتقي فيها
العقربان، وتقسيم يوم كامل إلى قطاعات}]
\label{pb:g6:lines:1}
الساعة منقلة تخبرك بالوقت: فوجهها دورة كاملة
قياسها $360^\circ$، تقسمها علامات الساعات الاثنتا عشرة إلى اثنتي عشرة
زاوية متساوية. وقد قاس \cref{exo:g6:lines:11} العقربين
عند 3 وعند 6؛ وتبني هذه المسألة النظرية الكاملة
--- بما في ذلك الأوقات التي \emph{لا} يشير فيها أيّ عقرب إلى علامة
--- وتجيب عن سؤال قلّ من الكبار من يصيبه («كم مرة يجلس
العقربان أحدهما فوق الآخر تمامًا؟»)، وتنتهي بتقسيم
يوم كامل إلى مخطط قطاعات.

\medskip
\noindent\textbf{الجزء الأول --- الكسور من الدورة.}
\begin{enumerate}
  \item الدورة الكاملة قياسها $360^\circ$. فكم درجة في
        نصف دورة، وفي ربع دورة، وفي جزء من اثني عشر من الدورة؟
  \item تقسم علامات الساعات الاثنتا عشرة وجه الساعة إلى اثنتي عشرة
        زاوية متساوية عند المركز. فكم درجة بين علامتين
        متجاورتين؟ اِستعد جواب
        \cref{exo:g6:lines:11} عند الساعة 3 بعدّ العلامات.
  \item أعطِ الزاوية بين العقربين عند الساعة 1، وعند
        الساعة 4، وعند الساعة 7. (عند الساعة 7 يفصل العقربان
        الوجه إلى زاويتين؛ و«الزاوية بين العقربين» تعني دائمًا
        \emph{الصغرى}.)
  \item يتمّ عقرب الدقائق دورة كاملة في $60$ دقيقة. فكم
        درجة يمسح في الدقيقة؟
  \item ينتقل عقرب الساعات من علامة إلى التي تليها ---
        $30^\circ$ --- في $60$ دقيقة. فكم درجة يمسح
        \emph{هو} في الدقيقة؟ (\omterm{def:g6:decimals:places}{عدد عشري}،
        \cref{ch:g6:decimals}.)
\end{enumerate}

\medskip
\noindent\textbf{الجزء الثاني --- أوقات لا يشير فيها شيء إلى
علامة.}
\begin{enumerate}[resume]
  \item عند 3:30 يشير عقرب الدقائق إلى العلامة $6$. فأين
        عقرب الساعات بالضبط؟ اِحسب زاوية كل عقرب انطلاقًا من
        العلامة $12$ (بالقياس في اتجاه عقارب الساعة)، واستنتج
        الزاوية بين العقربين.
  \item عند 6:30 يظنّ كثيرون أن العقربين أحدهما فوق
        الآخر. خمّن أولًا، ثم اِحسب الزاوية كما في السؤال
        6. فمن كان على صواب؟
  \item اِحسب الزاوية بين العقربين عند 9:15.
  \item اِشرح لماذا يجب، في مكان ما بين 1:00 و1:10، أن يكون
        العقربان أحدهما فوق الآخر تمامًا: أين
        عقرب الدقائق بالنسبة إلى عقرب الساعات عند 1:00، وأين
        هو عند 1:10؟ وأيّ عقرب يجري أسرع، ولماذا يحسم ذلك
        الأمر؟
  \item في اثنتي عشرة ساعة، كم مرة يجلس العقربان
        أحدهما فوق الآخر تمامًا؟ (يقع لقاء واحد في
        كل امتداد بين ساعتين متعاقبتين --- مع استثناء
        واحد: ماذا يحدث بين 11 و12؟ اُذكر أوقات اللقاء
        التقريبية وعُدَّها.) وكم مرة في
        يوم كامل؟
\end{enumerate}

\medskip
\noindent\textbf{الجزء الثالث --- الزوايا حول نقطة: مخططات
القطاعات.}
\begin{enumerate}[resume]
  \item تملأ قطاعات وجه الساعة الاثنا عشر الوجه
        كله: \omterm{def:g1:addition:def}{فمجموع} زواياها $360^\circ$. اِشرح لماذا
        يصحّ هذا في \emph{أيّ} مجموعة زوايا
        تشترك في رأس وتملأ دورة كاملة حوله، دون
        تراكب ودون فراغ.
  \item تسجّل زينة يومها ذا $24$ ساعة: نوم $9$ ساعات، ومدرسة
        $6$ ساعات، ولعب $3$ ساعات، ووجبات $2$ ساعة، وكل ما تبقّى $4$ ساعات.
        وهي تريد \emph{مخطط قطاعات}: قرصًا يأخذ فيه كل نشاط
        قطاعًا، بزوايا متناسبة مع الأزمنة.
        فأيّ \omterm{def:g6:fractions:def}{كسر} من اليوم هو كل نشاط، وكم
        درجة يأخذ قطاعه؟ وتحقّق من أن \omterm{def:g1:addition:def}{مجموع} الزوايا الخمس
        $360^\circ$.
  \item صِف، خطوة خطوة، كيف يُرسم مخطط قطاعات زينة
        بالفرجار والمنقلة
        (\cref{met:g6:lines:protractor}): أين يقع \omterm{def:g3:shapes:shapes}{نصف القطر}
        الأول، وكيف يبدأ كل قطاع جديد حيث ينتهي
        السابق.
  \item في مخطط قطاعات آخر ليوم، قطاع قياسه
        $90^\circ$ بالضبط. فأيّ \omterm{def:g6:fractions:def}{كسر} من القرص هو ذلك،
        وكم ساعة يمثّل؟
  \item الخاتمة، عودًا إلى الساعة: في $24$ ساعة، كم
        دورة كاملة يتمّ عقرب الساعات؟ وعقرب الدقائق؟ وعقرب
        الثواني؟ (أحد هذه الأجوبة يتجاوز الألف.)
\end{enumerate}
\end{problem}
